\documentclass[paperwidth=120cm,paperheight=150cm,portrait,fontscale=.22,margin=5cm,lmargin=.2cm,rmargin=.2cm]{baposter}

\usepackage{minted}
\definecolor{mintedbg}{rgb}{0.95,0.95,0.95}

\usepackage{upquote}
\AtBeginDocument{%
\def\PYZsq{\textquotesingle}%
}

\definecolor{bordercol}{RGB}{40,40,40}
\definecolor{headercol1}{RGB}{235,218,229}
\definecolor{headercol2}{RGB}{80,80,80}
\definecolor{headerfontcol}{RGB}{0,0,0}
\definecolor{boxcolor}{RGB}{255,235,205}

\renewcommand{\familydefault}{\sfdefault}

\def\mireval{\texttt{mir\char`_eval}}

\usepackage{tabularx}
\usepackage{booktabs}

\begin{document}

\begin{poster}{
  background=plain,
  bgColorOne=white,
  headerheight=0.1\textheight,
  eyecatcher=true,
  columns=2,
  borderColor=bordercol,
  headerColorOne=headercol1,
  headerColorTwo=headercol1,
  headerFontColor=headerfontcol,
  headerfont=\textmd,
  boxColorOne=boxcolor,
  linewidth=0.5pt,
  borderColor=black,
  headerborder=none,
  textborder=none,
}
{\includegraphics[height=8em]{labrosa-new-nobg.pdf}}
{\textmd{MIR\char`_EVAL:$\;$A$\;$TRANPARENT$\;$IMPLEMENTATION\\OF$\;$COMMON$\;$MIR$\;$METRICS}}
{Colin Raffel, Brian McFee, Eric J. Humphrey, Justin Salamon, Oriol Nieto, Dawen Liang,\\[.2em] and Daniel P. W. Ellis\\[.2em]}
{\includegraphics[height=6em]{marl-logo.pdf}}

\begin{posterbox}[column=0]{Abstract}
We present \mireval{}, an open source software library which provides a transparent and easy-to-use implementation of the most common metrics used to measure the performance of MIR algorithms.
We performed a quantitative comparison of \mireval{} to existing evaluation systems to explore differences in implementation.
\end{posterbox}

\begin{posterbox}[column=0,below=auto]{Design}

\mireval{} is a Python library which currently includes metrics for the following tasks: \textbf{Beat detection, chord estimation, pattern discovery, structural segmentation, melody extraction}, and \textbf{onset detection}.
Each task is given its own submodule, and each metric is defined as a separate function in each submodule.
Each task submodule also includes common data pre-processing steps for the task.
Every metric function includes detailed documentation, example usage, input validation, and references to the original paper which defined the metric.
\mireval{} also includes a submodule \texttt{io} which provides convenience functions for loading in task-specific data from common file formats.
In order to simplify the usage of \mireval{}, it is packaged with a set of ``evaluator'' scripts, one for each task.
These scripts include all code necessary to load in data, pre-process it, and compute all metrics for a given task.
The evaluators allow for \mireval{} to be called directly from the command line so that no knowledge of Python is necessary.

\end{posterbox}

\begin{posterbox}[column=1]{Comparison to Existing Implementations}


In order to validate the design choices made in \mireval{}, we compared the scores it produces to those reported by the evaluation systems used in MIREX.
Beyond pinpointing intentional differences in implementation, this process can also help find and fix bugs in either \mireval{} or the system it is being compared to.
For each task covered by \mireval{}, we obtained a collection of reference and estimated annotations and computed or obtained a score for each metric using \mireval{} and the evaluation system being compared to.
Then, for each reported score, we computed the relative change between the scores as their absolute difference divided by their mean.
Finally, we computed the average relative change across all examples in the obtained dataset for each score.
The number of algorithms, examples, and total number of scores for all tasks are summarized in the following table.

\begin{center}
\begin{tabular}{l r r r}
\toprule
Task & Algorithms & Examples & Scores\\
\midrule
Beat Detection & 20 & 679 & 13580\\
Segmentation & 8 & 1397 & 11176 \\
Onset Detection & 11 & 85 & 935 \\
Chord Estimation & 12 & 217 & 2604 \\
Melody & 1 & 20 & 20 \\
Pattern Discovery & 4 & 5 & 20 \\
\bottomrule
\end{tabular}
\end{center}

\vspace{-2.5mm}

\end{posterbox}

\begin{posterbox}[column=0,below=auto,span=2]{Average Relative Difference of \mireval{} vs. MIREX}

\begin{center}
\begin{tabularx}{.95\linewidth}{*{9}{>{\centering\arraybackslash}X}}
%\toprule
 \multicolumn{9}{ c }{\textbf{Beat Detection}} \\
  \cmidrule(lr){1-9}
F-measure  & Cemgil     & Goto       &  P-score   & CMLc       & CMLt       & AMLc       & AMLt       & In. Gain   \\
   0.703\% &    0.035\% &    0.054\% &    0.877\% &    0.161\% &    0.143\% &    0.137\% &    0.139\% &    9.174\% \\
  \midrule
\end{tabularx}
\begin{tabularx}{.95\linewidth}{*{9}{>{\centering\arraybackslash}X}}
 \multicolumn{9}{ c }{\textbf{Structural Segmentation}} \\
  \cmidrule(lr){1-9}
  NCE-Over   & NCE-under  & Pairwise F & Pairwise P & Pairwise R & Rand       & F@.5       & P@.5       & R@.5  \\
3.182\% &   11.082\% &    0.937\% &    0.942\% &    0.785\% &    0.291\% &    0.429\% &    0.088\% &    1.021\%  \\
  \midrule
\end{tabularx}
\begin{tabularx}{.95\linewidth}{*{8}{>{\centering\arraybackslash}X}}
 \multicolumn{5}{ c }{\textbf{Structural Segmentation (continued)}} & \multicolumn{3}{ c }{\textbf{Onset Detection}}\\
  \cmidrule(lr){1-5}
  \cmidrule(lr){6-8}
  F@3 &   P@3        & R@3 & Ref-est dev. & Est-ref dev. & F-measure  & Precision  & Recall     \\
 0.393\%  & 0.094\% &    0.954\% & 0.935\% &    0.000\% &  0.165\% &    0.165\% &    0.165\% \\
  \midrule
\end{tabularx}
\begin{tabularx}{.95\linewidth}{*{2}{>{\centering\arraybackslash}X} c *{1}{>{\centering\arraybackslash}X} c *{1}{>{\centering\arraybackslash}X} c *{1}{>{\centering\arraybackslash}X} c c}
 \multicolumn{5}{ c }{\textbf{Chord Estimation}} & \multicolumn{5}{ c }{\textbf{Melody Extraction}}\\
  \cmidrule(lr){1-5}
  \cmidrule(lr){6-10}
  Root       & Maj/min       & Maj/min + Inv & 7ths & 7ths + Inv & Overall & Raw pitch & Chroma & Voicing R & Voicing FA \\
 0.007\% & 0.163\% & 1.005\% & 0.483\% & 0.899\% & 0.070\% &    0.087\% &    0.114\% &    0.000\% &   10.095\% \\
%\bottomrule
\end{tabularx}
\end{center}

\end{posterbox}

\begin{posterbox}[column=0,below=auto,span=2,boxColorOne=white]{}

\begin{minted}[bgcolor=mintedbg, frame=single, framerule=0pt, fontsize=\large]{python}
import mir_eval
# Load in beat annotations
reference_beats = mir_eval.io.load_events('reference_beats.txt')
estimated_beats = mir_eval.io.load_events('estimated_beats.txt')
# scores will be a dictionary where the key is the metric name and the value is the score achieved
scores = mir_eval.beat.evaluate(reference_beats, estimated_beats)
# evaluate() will pass keyword args to the metric functions
scores = mir_eval.beat.evaluate(reference_beats, estimated_beats, f_measure_threshold=0.05)
# You can also perform pre-processing and compute each metric manually
reference_beats = mir_eval.beat.trim_beats(reference_beats)
estimated_beats = mir_eval.beat.trim_beats(estimated_beats)
f_score = mir_eval.beat.f_measure(reference_beats, estimated_beats, f_measure_threshold=0.05)
\end{minted}

$\;$\\[-1em]

\begin{minted}[bgcolor=mintedbg, frame=single, framerule=0pt, fontsize=\large]{c}
> ./onset_eval.py reference_onsets.txt estimated_onsets.txt -o output.json
reference_onsets.txt vs. estimated_onsets.txt
  F-measure : 0.531
  Precision : 0.481
     Recall : 0.591
Saving results to:  output.json
\end{minted}


\centering\huge
\texttt{http://www.github.com/craffel/mir\_eval}

\end{posterbox}

\end{poster}

\end{document}
